\section{Dữ liệu}

\begin{frame}[shrink=20]{EDA phát hiện gì trước khi chạy thí nghiệm}
  \small
  \begin{tabularx}{\textwidth}{@{}p{4.2cm} L@{}}
    \toprule
    \textbf{Phát hiện} & \textbf{Hệ quả lên thiết kế} \\
    \midrule
    Nhãn nhiễu \NumNoiseFloor{}--\NumNoiseCeiling{} &
      Chênh lệch dưới \NumNoiseFloor{} không được coi là thắng \\
    Trung vị \NumDistractorMedian{} đoạn đối thủ mỗi câu hỏi &
      Reranker là bắt buộc, không phải tùy chọn \\
    Câu hỏi neo vào tên riêng, ngày, số &
      Dự đoán sparse thắng dense --- \emph{trước} khi đo \\
    \NumCorpusRestored{} bài bị cắt đuôi &
      Dựng lại kho: \NumCorpusChunksVone{} $\rightarrow$ \NumCorpusChunks{} đoạn \\
    Notebook cũ đặt judge $=$ generator &
      Tách judge sang nhà cung cấp khác \\
    \bottomrule
  \end{tabularx}

  \takeaway{EDA không phải bước trang trí: nó viết sẵn giả thuyết cho từng vòng thí nghiệm.}
\end{frame}

\begin{frame}[shrink=20]{Hai cách đọc cùng một tập câu hỏi}{\texttt{original} và \texttt{resolved}}
  \small
  Nhiều câu hỏi NewsQA không đứng độc lập được --- chúng viết như đang nói tiếp
  một cuộc hội thoại:

  \begin{center}
    \begin{tabular}{@{}l l@{}}
      \texttt{original} & \emph{``How many were killed?''} \\
      \texttt{resolved} & \emph{``How many people were killed in the Kandahar bombing?''} \\
    \end{tabular}
  \end{center}

  \begin{itemize}
    \item Tỉ lệ câu phụ thuộc ngữ cảnh giảm từ \textbf{57,5\%} xuống \textbf{11,3\%} sau khi resolve.
    \item Bản \texttt{original} \textbf{thiên vị sparse}: câu cụt thì khớp từ vựng
          càng có lợi thế giả tạo. Đo trên đó là tự khen mình.
    \item Hệ thống thật luôn nhận câu hỏi đầy đủ ngữ cảnh $\rightarrow$
          \texttt{resolved} là tập \textbf{sát thực tế triển khai}.
  \end{itemize}

  \takeaway{Mọi số chính báo cáo theo \texttt{resolved}; \texttt{original} chỉ dùng làm cận dưới.}
\end{frame}
